\documentclass[../../main.tex]{subfiles}

\begin{document}

\chapter[Analisis Semantik]{Analisis Semantik: SDD, Translation Scheme,\\ dan Kode Antara}
\label{chap:semantik}

\begin{outcomebox}
\textbf{Sub-CPMK:} \subcpmk{1.5}\\[3pt]
Setelah menyelesaikan bab ini, mahasiswa mampu:
\begin{enumerate}[label=\arabic*.,leftmargin=*]
  \item menjelaskan ruang lingkup analisis semantik;
  \item membedakan atribut synthesized dan inherited;
  \item membuat tabel aturan semantik berbasis syntax-directed
        definition untuk ekspresi dan pernyataan subset bahasa C;
  \item menguji keluaran tabel aturan semantik pada pohon sintaks;
  \item menghasilkan keluaran postfix atau three-address code dari
        aturan semantik.
\end{enumerate}
\end{outcomebox}

\section*{Materi Pokok}
\begin{itemize}[leftmargin=*]
  \item Tugas dan ruang lingkup analisis semantik.
  \item Tabel aturan semantik (SDD) untuk ekspresi dan pernyataan.
  \item Translation scheme dan aksi semantik.
  \item Atribut synthesized dan inherited sebagai notasi pendukung.
  \item Kode antara: postfix, n-tuple, dan three-address code.
\end{itemize}

\subfile{section-01}
\subfile{section-03}
\subfile{section-04}
\subfile{section-02}
\subfile{section-05}

\section*{Contoh}
\addcontentsline{toc}{section}{Contoh}

\begin{contoh}[Atribut synthesized]
Untuk ekspresi \code{2 + 3.5}, tipe anak kiri adalah \code{int} dan
anak kanan \code{float}. Aturan \code{arith(int, float)}
menghasilkan \code{float}. Atribut tipe simpul akar dihitung dari
tipe anak, sehingga merupakan atribut synthesized.
\end{contoh}

\begin{contoh}[Atribut inherited]
Pada blok
\begin{lstlisting}[language=C]
{
    int x = 1;
    {
        int y = x + 1;
    }
}
\end{lstlisting}
lingkungan scope dalam mewarisi tabel simbol scope luar. Atribut
\code{env} pada blok dalam dihitung dari blok luar. Ini adalah
atribut inherited karena nilainya berasal dari orang tua.
\end{contoh}

\begin{contoh}[SDD ke translation scheme]
Aturan SDD
\[
E \to E_1 + E_2 \quad \{ E.type = \text{arith}(E_1.type,
E_2.type) \}
\]
dapat ditulis sebagai aksi di akhir produksi:
\begin{lstlisting}[language=Bison]
expr : expr T_PLUS expr
     { $$.type = type_arith($1.type, $3.type); }
     ;
\end{lstlisting}
Aksi dijalankan saat reduksi, setelah tipe kedua operand diketahui.
\end{contoh}

\begin{contoh}[Dari SDD ke postfix]
Untuk aturan
\[
E \to E_1 + T \quad \{\, E.code = E_1.code \,\|\, T.code \,\|\,
\text{``+''} \,\}
\]
ekspresi \code{(2 + 3) * 4} menghasilkan postfix
\code{2 3 + 4 *}. Atribut \code{code} merupakan atribut synthesized:
nilainya dihitung dari kode anak-anaknya.
\end{contoh}

\begin{contoh}[Dari SDD ke three-address code]
Untuk pernyataan \code{a = b + c * d}, aturan semantik dapat
membangkitkan instruksi
\[
t_1 = c * d, \qquad t_2 = b + t_1, \qquad a = t_2.
\]
Setiap hasil antara disimpan pada nama sementara. Urutan instruksi
mengikuti dependensi pohon: perkalian diselesaikan sebelum
penjumlahan, dan penjumlahan sebelum assignment.
\end{contoh}

\begin{contoh}[Aturan pernyataan pada pohon sintaks]
Untuk program
\begin{lstlisting}[language=C]
int a = 2;
float b = 3.5;
a = b + 1;
\end{lstlisting}
pohon sintaks memuat deklarasi \(a\), deklarasi \(b\), ekspresi
\(b + 1\), dan assignment \(a = b + 1\). Aturan
\(\text{declare}(a,\code{int})\) dan
\(\text{declare}(b,\code{float})\) mengisi tabel simbol. Aturan
\(\text{arith}(\code{float},\code{int})\) menghasilkan
\code{float}. Aturan
\(\text{check\_assign}(\code{int},\code{float})\) menghasilkan
\code{error} karena assignment memerlukan cast.
\end{contoh}

\section*{Aktivitas Pembelajaran dan Praktikum}
\addcontentsline{toc}{section}{Aktivitas Pembelajaran dan Praktikum}

\begin{table}[H]
\centering
\caption{Rencana aktivitas Bab XIV}
\label{tab:ch14-aktivitas}
\begin{tabularx}{\textwidth}{@{}llX@{}}
\toprule
\textbf{Sesi} & \textbf{Bentuk} & \textbf{Kegiatan} \\
\midrule
1 & Ceramah & Ruang lingkup analisis semantik \\
2 & Latihan & Menulis tabel aturan semantik ekspresi dan pernyataan \\
3 & Diskusi & Pengujian aturan pada pohon sintaks \\
4 & Praktikum & P4-A: tabel aturan, uji pohon, dan inferensi tipe \\
\bottomrule
\end{tabularx}
\end{table}

\subsection*{P4-A: Tabel Aturan Semantik dan Inferensi Tipe}
\textbf{Tujuan.} Membuat tabel aturan semantik untuk ekspresi dan
pernyataan, mengujinya pada pohon sintaks, lalu
mengimplementasikan inferensi tipe.

\textbf{Langkah kerja.}
\begin{enumerate}[label=\arabic*.,leftmargin=*]
  \item Lengkapi tabel aturan semantik ekspresi pada
        \tabref{tab:ch14-sdd} dan tabel pernyataan pada
        \tabref{tab:ch14-sdd-stmt}.
  \item Pilih satu program yang memuat deklarasi, assignment, dan
        percabangan atau perulangan; gambar pohon sintaksnya.
  \item Terapkan setiap aturan pada simpul pohon dan catat hasil atau
        \code{TYPE\_ERROR}.
  \item Lengkapi \file{include/types.h} dan \file{src/types.c}.
  \item Implementasikan \code{infer\_expr} dan \code{check\_stmt}.
  \item Uji minimal 15 ekspresi dan 5 pernyataan, termasuk campuran
        tipe dan identifier yang tidak dideklarasikan.
  \item Kumpulkan tabel, pohon teranotasi, dan keluaran pengujian.
\end{enumerate}

\textbf{Definisi selesai.}
\begin{checklist}
  \item Tabel aturan ekspresi dan pernyataan lengkap.
  \item Pohon sintaks teranotasi dengan tipe atau error.
  \item Minimal 15 ekspresi dan 5 pernyataan diuji.
  \item \code{TYPE\_ERROR} dipropagasikan tanpa pesan berulang.
  \item Laporan dan keluaran pengujian diserahkan.
\end{checklist}

P4-A menjadi bagian Praktikum 4 (12\%) yang dilaporkan pada Bab XV.

\section*{Latihan dan Refleksi}
\addcontentsline{toc}{section}{Latihan dan Refleksi}

\subsection*{Latihan}
\begin{enumerate}[label=\arabic*.,leftmargin=*]
  \item Jelaskan perbedaan analisis sintaksis dan semantik.
  \item Definisikan atribut synthesized dan inherited.
  \item Tuliskan tabel SDD untuk tipe ekspresi penjumlahan dan
        perkalian, serta untuk assignment, if, dan while.
  \item Jelaskan perbedaan SDD dan translation scheme.
  \item Berikan contoh atribut inherited pada scope.
  \item Tuliskan postfix untuk \code{(a + b) * (c - d)}.
  \item Tuliskan three-address code untuk \code{x = a * b + c}.
  \item Mengapa analisis semantik memerlukan tabel simbol?
  \item Uji tabel aturan pada pohon sintaks untuk program
        \code{int a; float b; a = b + 1;} dan catat hasil atau error
        setiap simpul.
\end{enumerate}

\subsection*{Refleksi}
\begin{reflectionbox}
\begin{enumerate}[label=\arabic*.,leftmargin=*]
  \item Apakah Anda dapat membedakan SDD dan translation scheme?
  \item Bagian mana dari pembentukan kode antara yang paling menantang?
  \item Mengapa urutan instruksi harus mengikuti dependensi pohon?
\end{enumerate}
\end{reflectionbox}

\section*{Asesmen}
\addcontentsline{toc}{section}{Asesmen}

Asesmen Bab XIV menjadi bagian Praktikum 4 (8\%) dan Tugas Teori
(8\%). Fokusnya adalah SDD, translation scheme, dan inferensi tipe.

\begin{table}[H]
\centering
\caption{Instrumen asesmen Bab XIV}
\label{tab:ch14-asesmen}
\begin{tabularx}{\textwidth}{@{}lXl@{}}
\toprule
\textbf{Instrumen} & \textbf{Indikator} & \textbf{Bobot} \\
\midrule
P4-A & Inferensi tipe ekspresi & bagian P4 \\
Tugas Teori & SDD dan translation scheme & 8\% \\
Kuis & Atribut synthesized/inherited & formatif \\
\bottomrule
\end{tabularx}
\end{table}

Kriteria kelulusan: mahasiswa dapat menulis SDD, menjelaskan
perbedaan synthesized dan inherited, serta mengimplementasikan
inferensi tipe dasar.

\section*{Checklist Pencapaian Kompetensi}
\addcontentsline{toc}{section}{Checklist Pencapaian Kompetensi}

\begin{checklist}
  \item Saya dapat menjelaskan ruang lingkup analisis semantik.
  \item Saya dapat membedakan atribut synthesized dan inherited.
  \item Saya dapat membuat tabel aturan semantik untuk ekspresi.
  \item Saya dapat membuat tabel aturan semantik untuk pernyataan.
  \item Saya dapat menguji aturan semantik pada pohon sintaks.
  \item Saya dapat menjelaskan translation scheme.
  \item Saya dapat menentukan tipe ekspresi dan pernyataan campuran.
  \item Saya dapat mempropagasikan \code{TYPE\_ERROR}.
  \item Saya dapat mengimplementasikan inferensi tipe dasar.
  \item Saya sudah menyelesaikan P4-A.
\end{checklist}

\section*{Rangkuman}
\addcontentsline{toc}{section}{Rangkuman}

Pada bab ini telah dibahas:
\begin{itemize}[leftmargin=*]
  \item tugas dan ruang lingkup analisis semantik;
  \item attribute grammar dengan atribut synthesized dan inherited;
  \item tabel aturan semantik untuk ekspresi dan pernyataan;
  \item pengujian aturan pada pohon sintaks;
  \item translation scheme dan penerapannya;
  \item kode antara: postfix, n-tuple, dan three-address code.
\end{itemize}

Bab XV akan membahas tabel simbol, manajemen scope, dan pemeriksaan
tipe secara lebih rinci. Implementasi tabel simbol menjadi fondasi
bagi analisis semantik yang lengkap.

\section*{Evaluasi}
\addcontentsline{toc}{section}{Evaluasi}

\begin{enumerate}[label=\arabic*.,leftmargin=*]
  \item Jelaskan tiga tugas analisis semantik.
  \item Definisikan atribut synthesized dan inherited.
  \item Tuliskan tabel SDD untuk tipe ekspresi dan pernyataan.
  \item Jelaskan perbedaan SDD dan translation scheme.
  \item Tuliskan postfix untuk \code{a + b * c - d}.
  \item Tuliskan quadruple dan three-address code untuk
        \code{x = a * b + c}.
  \item Mengapa tabel simbol diperlukan?
  \item Berikan contoh kesalahan semantik yang tidak dapat dideteksi
        parser.
  \item Uji tabel aturan pada pohon sintaks untuk program
        \code{int a; float b; a = b + 1;} dan catat hasil atau error
        setiap simpul.
\end{enumerate}

\subsection*{Kunci dan Rubrik Singkat}
\begin{itemize}[leftmargin=*]
  \item Soal 1--4: ketepatan konsep (70\%), contoh (30\%).
  \item Soal 5--6: kebenaran postfix dan three-address code (70\%),
        penjelasan (30\%).
  \item Soal 7--9: analisis (60\%), argumentasi (40\%).
\end{itemize}


\ifSubfilesClassLoaded{%
  \bibliographystyle{apalike}%
  \bibliography{../../references}%
}{}

\end{document}
